\chapter*{증거와 표현의 규칙}
\addcontentsline{toc}{chapter}{증거와 표현의 규칙}

\begin{factbox}
``활성값에 정보가 존재한다'', ``모델이 그 정보를 사용한다'', ``특정 구성요소가 그 정보를 계산한다''는 서로 다른 주장이다.
첫째는 decodability, 둘째는 causal relevance, 셋째는 mechanism에 관한 주장이다.
한 종류의 증거만으로 나머지를 자동으로 결론낼 수 없다.
\end{factbox}

\section*{네 가지 서술 등급}

본문의 색상 상자는 주장 종류를 표시한다.

\begin{description}[style=nextline,leftmargin=2em]
  \item[수학적 사실 또는 명시된 구현]
  정의에서 유도되는 등식, 증명된 명제, 또는 특정 공개 코드가 실제로 수행하는 연산이다.
  구현에 관한 문장은 모델·버전·코드 경로를 제한한다.

  \item[실험 결과]
  인용한 연구가 정한 모델, 데이터, metric과 intervention 아래에서 관찰한 결과다.
  표본 밖의 모델이나 과제로 일반화하지 않는다.

  \item[제한된 해석]
  관찰을 압축하는 유용한 모델이나 연구자들의 해석이지만 유일한 설명으로 증명되지 않은 내용이다.
  ``attention은 정보를 가져온다'', ``MLP는 memory다'' 같은 표현은 이 등급에서만 사용한다.

  \item[한계 또는 열린 문제]
  현재 방법으로 식별되지 않거나 논문 간 결론이 충돌하는 부분이다.
  빈칸을 그럴듯한 이야기로 채우지 않는다.
\end{description}

\section*{설명의 다섯 수준}

\begin{table}[htbp]
\centering
\caption{LLM을 설명한다는 말의 서로 다른 수준}
\label{tab:levels-of-explanation}
\begin{tabularx}{\textwidth}{@{}p{0.17\textwidth}YY@{}}
\toprule
수준 & 질문 & 검증 기준 \\
\midrule
구현 & 어떤 연산을 실행하는가? & 코드와 텐서 값을 재현한다. \\
수학 & 입력에서 출력까지 함수는 무엇인가? & 수식의 차원과 등식이 맞는다. \\
표현 & 어떤 변수가 내부 상태에서 복원되는가? & held-out data와 적절한 control을 사용한다. \\
인과 & 어떤 구성요소가 행동에 영향을 주는가? & 명시된 intervention과 metric으로 효과를 재현한다. \\
알고리즘 & 구성요소가 어떤 절차로 협력하는가? & 회로가 행동을 충실히 재현하고 반사실적 예측을 맞힌다. \\
\bottomrule
\end{tabularx}
\end{table}

완전한 설명은 최소한 이 수준들을 연결해야 한다.
그러나 세부 수치 계산을 그대로 나열하면 완전하지만 이해하기 어렵고, 짧은 의사코드는 이해하기 쉽지만 실제 계산을 놓칠 수 있다.
따라서 ``완전성''과 ``압축성''은 별개의 평가축이다.

\section*{경험적 주장에 필요한 여섯 항목}

경험적 주장은 가능한 경우 다음 여섯 정보를 같은 문단이나 인접 표에 기록한다.

\begin{enumerate}
  \item 모델 이름, 크기, checkpoint와 architecture;
  \item 입력 분포와 표본 선택 규칙;
  \item 설명하려는 behavior와 평가 metric;
  \item 관찰하거나 개입한 activation, edge, weight 또는 feature;
  \item clean, corrupted, patched 조건의 정확한 차이;
  \item 결론이 적용되는 범위와 알려진 대안 설명.
\end{enumerate}

특히 ``중요하다''는 말은 metric을 생략하면 불완전하다.
정답 logit, logit difference, probability, KL divergence, cross-entropy, accuracy는 서로 다른 양이며 intervention의 순위도 바꿀 수 있다 \citep{zhang2024activation,miller2024faithfulness}.

\section*{이 책에서 금지하는 추론}

\begin{itemize}
  \item attention weight가 크다는 사실만으로 그 source token이 출력의 원인이라고 쓰지 않는다.
  \item linear probe가 성공했다는 사실만으로 원 모델이 그 방향을 사용한다고 쓰지 않는다.
  \item 한 위치를 편집할 수 있다는 사실만으로 지식이 오직 그곳에 저장됐다고 쓰지 않는다.
  \item ablation 효과가 작다는 사실만으로 구성요소가 불필요하다고 쓰지 않는다.
    downstream self-repair가 효과를 보상할 수 있다.
  \item 생성된 chain-of-thought를 내부 계산의 충실한 transcript라고 가정하지 않는다.
  \item 소형·합성 과제의 회로를 자연어 전반의 보편 알고리즘이라고 쓰지 않는다.
\end{itemize}

이 주의들은 단순한 철학적 경고가 아니다.
probing control \citep{hewitt2019probes}, attention 설명 논쟁 \citep{jain2019attention,wiegreffe2019attention}, activation patching의 방법 민감도 \citep{zhang2024activation}, self-repair \citep{mcgrath2023hydra}, localization과 editing의 불일치 \citep{hase2023localization}, CoT 비충실성 \citep{turpin2023unfaithful}에서 각각 실험적으로 문제가 드러났다.

\section*{출처 정책}

핵심 주장은 원 논문, 공식 모델 보고서, 공개 구현 또는 표준 교재를 우선한다.
survey는 문헌 지도를 찾는 데 사용할 수 있지만 원 결과의 대체 인용으로 쓰지 않는다.
기업 연구팀이 자기 모델을 분석한 기술 보고서는 primary evidence이지만 독립 재현이 아니라는 점을 유지한다.
arXiv preprint는 동료심사가 완료된 논문과 동일하게 표시하지 않는다.

\begin{limitbox}
어떤 문헌 정책도 오류 가능성을 0으로 만들지는 않는다.
정확성은 출처의 권위만으로 보장되지 않으며, 수식 검산, 코드 대조, 재현, 반대 결과 확인과 판본 갱신이 필요하다.
\end{limitbox}
